\documentclass[../../../include/open-logic-section]{subfiles}

\begin{document}

\olfileid[hi]{sth}{spine}{Valphabasic}
\olsection{चरणों के मूल गुण}

$V_\alpha$ की परिभाषा का आधारभूत महत्त्व स्पष्ट करने के लिए हम उनके विषय
में कुछ मूल परिणाम प्रस्तुत करेंगे। एक परिभाषा से आरंभ करें:\footnote{
``सामर्थ्यवान'' (potent) के लिए कोई मानक शब्दावली नहीं; यह
\citet{ButtonLT1} द्वारा प्रयुक्त नाम का अनुवाद है।}
\begin{defn}
	समुच्चय $A$ \emph{सामर्थ्यवान} तभी और केवल तभी है जब
	$\forall x((\exists y \in A)x \subseteq y \lif x \in A)$।
\end{defn}
\begin{lem}\ollabel{Valphabasicprops}
प्रत्येक क्रमसूचक $\alpha$ के लिए:
\begin{enumerate}
	\item\ollabel{Valphatrans} प्रत्येक $V_\alpha$ संक्रमणशील है।
	\item\ollabel{Valphapotent} प्रत्येक $V_\alpha$ सामर्थ्यवान है।
	\item\ollabel{Valphacum} यदि $\gamma \in \alpha$, तो
	$V_\gamma \in V_\alpha$ (और इसलिए \olref{Valphatrans} से
	$V_\gamma \subseteq V_\alpha$ भी)।
\end{enumerate}
\end{lem}

\begin{proof}
हम इसे (युगपत) परिमितेतर आगमन से सिद्ध करते हैं। आगमन के लिए मान लें कि
प्रत्येक क्रमसूचक $\beta < \alpha$ के लिए
\olref{Valphatrans}--\olref{Valphacum} लागू होते हैं।

$\alpha = \emptyset$ का मामला तुच्छ है।

मान लें $\alpha = \ordsucc{\beta}$। \olref{Valphacum} दिखाने के लिए, यदि
$\gamma \in \alpha$, तो परिकल्पना से $V_\gamma \subseteq V_\beta$, अतः
$V_\gamma \in \Pow{V_\beta} = V_\alpha$। \olref{Valphapotent} दिखाने के
लिए मान लें $A \subseteq B \in V_\alpha$, अर्थात
$A \subseteq B \subseteq V_\beta$; तब $A \subseteq V_\beta$, इसलिए
$A \in V_\alpha$। \olref{Valphatrans} दिखाने के लिए ध्यान दें कि यदि
$x \in A \in V_\alpha$, तो $A \subseteq V_\beta$, अतः
$x \in V_\beta$; परिकल्पना से $V_\beta$ संक्रमणशील है, इसलिए
$x \subseteq V_\beta$ और इस प्रकार $x \in V_\alpha$।

मान लें $\alpha$ सीमा क्रमसूचक है। \olref{Valphacum} दिखाने के लिए, यदि
$\gamma \in \alpha$, तो $\gamma \in \ordsucc{\gamma} \in \alpha$;
इसलिए मान्यता से $V_\gamma \in V_{\ordsucc{\gamma}}$, और फलतः
$V_\gamma \in \bigcup_{\beta \in \alpha} V_\beta = V_\alpha$।
\olref{Valphatrans} और \olref{Valphapotent} दिखाने के लिए बस देखें कि
संक्रमणशील (क्रमशः सामर्थ्यवान) समुच्चयों का सम्मिलन संक्रमणशील (क्रमशः
सामर्थ्यवान) होता है।
\end{proof}

\begin{lem}\ollabel{Valphanotref}
प्रत्येक क्रमसूचक $\alpha$ के लिए $V_\alpha \notin V_\alpha$।
\end{lem}

\begin{proof}
परिमितेतर आगमन से। स्पष्टतः $V_\emptyset \notin V_\emptyset$।

यदि $V_{\ordsucc{\alpha}} \in V_{\ordsucc{\alpha}} = \Pow{V_\alpha}$,
तो $V_{\ordsucc{\alpha}} \subseteq V_\alpha$; और चूँकि
\olref{Valphabasicprops} से $V_\alpha \in V_{\ordsucc{\alpha}}$, इसलिए
$V_\alpha \in V_\alpha$। इसके विपरीत: यदि
$V_\alpha \notin V_\alpha$, तो
$V_{\ordsucc{\alpha}} \notin V_{\ordsucc{\alpha}}$।

यदि $\alpha$ सीमा है और
$V_\alpha \in V_\alpha = \bigcup_{\beta \in \alpha}V_\beta$, तो किसी
$\beta \in \alpha$ के लिए $V_\alpha \in V_\beta$; किंतु तब
$V_\beta \in V_\alpha$ भी, जिससे \olref{Valphabasicprops} (दो बार) द्वारा
$V_\beta \in V_\beta$। इसके विपरीत यदि सभी $\beta \in \alpha$ के लिए
$V_\beta \notin V_\beta$, तो $V_\alpha \notin V_\alpha$।
\end{proof}

\begin{cor}
किन्हीं क्रमसूचकों $\alpha, \beta$ के लिए: $\alpha \in \beta$ तभी और
केवल तभी जब $V_\alpha \in V_\beta$।
\end{cor}

\begin{proof}
\Olref{Valphabasicprops} एक दिशा देता है। इसके विपरीत मान लें
$V_\alpha \in V_\beta$। तब \olref{Valphanotref} से
$\alpha \neq \beta$; और $\beta \notin \alpha$, क्योंकि अन्यथा
$V_\beta \in V_\alpha$ और फिर \olref{Valphabasicprops} (दो बार) से
$V_\beta \in V_\beta$, जो \olref{Valphanotref} के विरुद्ध है। अतः त्रिकोटी
से $\alpha \in \beta$।
\end{proof}
\noindent
यह सब हमें प्रत्येक $V_\alpha$ को पदानुक्रम का $\alpha$वाँ चरण मानने देता
है। कारण यह है।

निश्चय ही हमारे $V_\alpha$ को \emph{पुनरावर्ती} प्रक्रिया में बना माना जा
सकता है, क्योंकि क्रमसूचकों का हमारा उपयोग पुनरावृत्ति की धारणा का अनुसरण
करता है। इसके अतिरिक्त यदि एक चरण दूसरे से पहले बना, अर्थात
$V_\beta \in V_\alpha$, अर्थात $\beta \in \alpha$, तो हमारी निर्माण
प्रक्रिया \emph{संचयी} है, क्योंकि $V_\beta \subseteq V_\alpha$। अंततः हम
वास्तव में पहले के किसी चरण पर उपलब्ध समुच्चयों के \emph{सभी} संभव संग्रह
बना रहे हैं, क्योंकि प्रत्येक उत्तरवर्ती चरण $V_{\ordsucc{\alpha}}$ अपने
पूर्ववर्ती $V_\alpha$ का घात समुच्चय है।

संक्षेप में: $\ZFminus$ के साथ समुच्चयों के संचयी-पुनरावर्ती पदानुक्रम की
अपनी दृष्टि को व्यक्त करने का काम \emph{लगभग} पूरा है। (यद्यपि निश्चय ही
प्रतिस्थापन का औचित्य देना अब भी शेष है।)

\end{document}
